\documentclass[11pt,a4paper]{article}

% Compilación: pdfLaTeX.
\usepackage[T1]{fontenc}
\usepackage[utf8]{inputenc}
\usepackage[spanish,es-nodecimaldot]{babel}
\spanishplainpercent
\usepackage{helvet}
\IfFileExists{tinos.sty}{\usepackage{tinos}}{\usepackage{mathptmx}}
\renewcommand{\familydefault}{\sfdefault}
\usepackage[
  a4paper,
  left=15.9mm,
  right=16.3mm,
  top=25mm,
  bottom=25mm,
  headheight=5mm,
  headsep=6mm,
  footskip=13mm
]{geometry}
\usepackage{amsmath,amssymb}
\usepackage{graphicx}
\graphicspath{{figures/}}
\usepackage{xcolor}
\usepackage{tikz}
\usetikzlibrary{arrows.meta,babel}
\usepackage{fancyhdr}
\usepackage{setspace}
\usepackage{microtype}
\usepackage{booktabs,tabularx,array}
\usepackage{hyperref}

\definecolor{UAHBlue}{RGB}{52,74,139}
\colorlet{uahblue}{UAHBlue}
\colorlet{headingblue}{UAHBlue}
\hypersetup{
  colorlinks=true,urlcolor=UAHBlue,linkcolor=UAHBlue,
  pdftitle={Práctica final de Sistemas de Control Inteligente},
  pdfauthor={Departamento de Electrónica. Universidad de Alcalá},
  pdfsubject={Control borroso y neuronal de un vehículo autónomo}
}

\setlength{\parindent}{0pt}
\setlength{\parskip}{1.5mm}
\setstretch{1.0}
\emergencystretch=1.5em
\renewcommand{\arraystretch}{1.15}

\pagestyle{fancy}
\fancyhf{}
\fancyhead[L]{\fontsize{8.5}{9.5}\selectfont\itshape Práctica final}
\fancyhead[R]{\fontsize{8.5}{9.5}\selectfont\itshape Sistemas de Control Inteligente}
\renewcommand{\headrulewidth}{0pt}
\fancyfoot[L]{%
  \begin{minipage}{\textwidth}
    \hrule height 0.35pt
    \vspace{1.2mm}%
    \fontsize{8.1}{9}\selectfont
    Departamento de Electrónica. Universidad de Alcalá \hfill \thepage
  \end{minipage}}
\renewcommand{\footrulewidth}{0pt}

\newcommand{\problem}[1]{%
  \par\addvspace{4.2mm}%
  {\rmfamily\fontsize{16}{18}\selectfont\color{headingblue}#1\par}%
  \nobreak\vspace{1.6mm}\nobreak
}
\newcommand{\romanitem}[2]{%
  \par\vspace{2mm}%
  \noindent\hangindent=8mm\hangafter=1%
  \makebox[8mm][r]{\textbf{#1.}\hspace{1.5mm}}#2\par
}
\newcommand{\centerEq}[1]{%
  \begin{center}$#1$\end{center}
}

\begin{document}
\thispagestyle{empty}

\begin{tikzpicture}[remember picture,overlay]
  \IfFileExists{figures/UAH-cabecera.png}{%
    \node[anchor=north west,inner sep=0pt]
      at ([xshift=-25mm,yshift=5mm]current page.north west)
      {\includegraphics[width=130mm]{figures/UAH-cabecera.png}};
  }{%
    \IfFileExists{UAH-cabecera.png}{%
      \node[anchor=north west,inner sep=0pt]
        at ([xshift=-25mm,yshift=5mm]current page.north west)
        {\includegraphics[width=130mm]{UAH-cabecera.png}};
    }{%
      \node[anchor=north west,text=uahblue,
            font=\bfseries\fontsize{19}{22}\selectfont]
        at ([xshift=20mm,yshift=-21mm]current page.north west)
        {Universidad de Alcalá};
    }
  }
  \fill[uahblue]
    ([yshift=-72.5mm]current page.north west)
    rectangle
    ([xshift=15.3mm,yshift=-97.5mm]current page.north west);
  \fill[uahblue]
    ([xshift=19.8mm,yshift=-72.5mm]current page.north west)
    rectangle
    ([xshift=207.3mm,yshift=-97.5mm]current page.north west);
  \node[text=white,font=\fontsize{29}{31}\selectfont,anchor=center]
    at ([xshift=73.5mm,yshift=-85mm]current page.north west)
    {Práctica final};
  \node[text=uahblue,font=\fontsize{35}{38}\selectfont,anchor=north]
    at ([xshift=105mm,yshift=-116.5mm]current page.north west)
    {Sistemas de Control Inteligente};
  \node[text=uahblue,font=\fontsize{20}{24}\selectfont,
        align=center,text width=175mm,anchor=north]
    at ([xshift=105mm,yshift=-143.5mm]current page.north west)
    {Control borroso y neuronal\\de un vehículo autónomo};
  \node[text=uahblue,font=\bfseries\fontsize{18}{20}\selectfont,anchor=north]
    at ([xshift=105mm,yshift=-180mm]current page.north west)
    {Grado en Matemáticas y Computación};
  \node[text=uahblue,font=\bfseries\fontsize{24}{26}\selectfont,anchor=north]
    at ([xshift=105mm,yshift=-203mm]current page.north west)
    {Universidad de Alcalá};
  \draw[uahblue,line width=0.45pt]
    ([xshift=83.5mm,yshift=-216mm]current page.north west)
    --
    ([xshift=126.5mm,yshift=-216mm]current page.north west);
\end{tikzpicture}
\null
\newpage
\setcounter{page}{2}
\fontsize{11}{13.1}\selectfont

\problem{1. Objetivos}

El objetivo de la práctica es diseñar y comparar dos controladores capaces
de conducir un vehículo autónomo por un circuito cerrado: un controlador
borroso de tipo \textbf{Mamdani}, basado en sonares, y un controlador
\textbf{neuronal}, basado en el LiDAR.

Ambos controladores deberán completar una vuelta en el menor tiempo posible,
mantener el vehículo en el carril derecho y adelantar al tráfico lento.
El carril contrario podrá utilizarse durante los adelantamientos, regresando
al carril derecho cuando exista espacio suficiente para hacerlo.
Se valorarán la ausencia de colisiones y la reducción del tiempo de ocupación
del carril contrario.

La conducción y la obtención de datos serán \textbf{automáticas}.
Los datos de entrenamiento de la red se obtendrán mediante ejecuciones del
controlador borroso desarrollado en la primera parte.

\problem{2. Plataforma robótica y circuito}

El entorno utiliza \textbf{MVSim y ROS~2}. El control se ejecuta en
\textbf{MATLAB/Simulink} y se comunica con el simulador mediante ROS Toolbox.
La instalación y el lanzamiento del simulador se describen en la
\href{https://robesafe-uah.github.io/sci_practica_final/}{guía del proyecto}.

El vehículo es un F1TENTH de tipo Ackermann. Solo se permite avanzar;
el vehículo necesita velocidad longitudinal para cambiar de orientación
y no puede girar sobre sí mismo.

\begin{center}
\begin{tabularx}{0.94\textwidth}{@{}X l@{}}
\toprule
\textbf{Parámetro} & \textbf{Valor}\\
\midrule
Velocidad de avance permitida & $0$ a $50\ \mathrm{km/h}$\\
Mando de dirección del controlador & $-45^\circ$ a $+45^\circ$\\
Distancia entre ejes del modelo & $0{,}3302\ \mathrm{m}$\\
Frecuencia de simulación y sensores & $20\ \mathrm{Hz}$\\
Periodo de muestreo del controlador & $T_s=0{,}05\ \mathrm{s}$\\
\bottomrule
\end{tabularx}
\end{center}

El circuito está formado por cuatro rectas unidas mediante cuatro curvas
circulares. Las dimensiones de su línea central son
\textbf{$40\times26\ \mathrm{m}$}, con radio de curva de
\textbf{$8\ \mathrm{m}$}. La carretera tiene una anchura total de
\textbf{$1{,}5\ \mathrm{m}$}, dividida en dos carriles de
\textbf{$0{,}75\ \mathrm{m}$}, uno por sentido.

El vehículo del alumno comienza en el carril derecho y recorre el circuito
en sentido antihorario. Los vehículos de tráfico son \textbf{naranjas},
circulan lentamente por el mismo carril y en el mismo sentido.
El escenario permite configurar entre \textbf{0 y 6 vehículos de tráfico}.

El coche del alumno y el tráfico permanecen detenidos hasta la primera
orden válida en \texttt{/r0/cmd\_vel}. Una vez iniciado, el tráfico
continúa circulando aunque el coche del alumno se detenga.

\newpage
\problem{3. Sensores e interfaz de control}

El vehículo dispone de un \textbf{Hokuyo simulado} con cobertura de
\textbf{$270^\circ$}, \textbf{1081 medidas} por barrido y alcance máximo
de \textbf{$30\ \mathrm{m}$}. El ángulo cero apunta hacia delante;
los ángulos positivos corresponden al lado izquierdo del vehículo.
Los $90^\circ$ posteriores quedan fuera de cobertura.

Los ocho sonares son \textbf{medidas virtuales derivadas del LiDAR}.
Cada uno publica la media de las distancias válidas de un sector de
$33{,}75^\circ$. La numeración avanza desde el lado derecho hacia el izquierdo.

\begin{figure}[!ht]
\centering
\begin{minipage}[c]{0.36\textwidth}
\centering
\begin{tabular}{@{}c r@{}}
\toprule
\textbf{Sonar} & \textbf{Ángulo central}\\
\midrule
0 & $-118{,}125^\circ$\\
1 & $-84{,}375^\circ$\\
2 & $-50{,}625^\circ$\\
3 & $-16{,}875^\circ$\\
4 & $+16{,}875^\circ$\\
5 & $+50{,}625^\circ$\\
6 & $+84{,}375^\circ$\\
7 & $+118{,}125^\circ$\\
\bottomrule
\end{tabular}
\end{minipage}
\hfill
\begin{minipage}[c]{0.58\textwidth}
\centering
\begin{tikzpicture}[scale=0.9,>=Stealth]
  \foreach \i in {0,...,7}{
    \pgfmathsetmacro{\a}{-45+33.75*\i}
    \pgfmathsetmacro{\b}{\a+33.75}
    \pgfmathsetmacro{\m}{(\a+\b)/2}
    \filldraw[fill=uahblue!8,draw=uahblue!55,line width=0.4pt]
      (0,0) -- (\a:2.35)
      arc[start angle=\a,end angle=\b,radius=2.35] -- cycle;
    \node[text=uahblue,font=\small\bfseries] at (\m:2.64) {\i};
  }
  \draw[fill=white,draw=uahblue,line width=0.8pt,rounded corners=2pt]
    (-0.35,-0.95) rectangle (0.35,0.40);
  \foreach \x in {-0.43,0.35}{
    \foreach \y in {-0.75,0.05}{
      \fill[black!75] (\x,\y) rectangle ++(0.08,0.28);
    }
  }
  \draw[->,uahblue,line width=0.9pt] (0,0.60) -- (0,1.40);
  \node[font=\small,text=uahblue] at (0,1.65) {Delante};
  \fill[uahblue] (0,0) circle[radius=1.7pt];
  \node[font=\scriptsize,text=uahblue] at (0,-0.32) {LiDAR};
  \node[font=\small,text=black!65,align=center] at (0,-2.15)
    {Zona sin cobertura\\$90^\circ$};
\end{tikzpicture}
\end{minipage}
\caption{Sonares virtuales y cobertura del LiDAR, en vista superior.
El dibujo representa la distribución angular, no el alcance a escala.}
\end{figure}

\begin{center}
\small
\begin{tabularx}{\textwidth}{@{}p{4.1cm} p{5.1cm} X@{}}
\toprule
\textbf{Topic} & \textbf{Tipo ROS~2} & \textbf{Contenido}\\
\midrule
\texttt{/r0/cmd\_vel} & \texttt{geometry\_msgs/msg/Twist}
  & Orden de velocidad y giro.\\
\texttt{/r0/scan} & \texttt{sensor\_msgs/msg/LaserScan}
  & Barrido LiDAR.\\
\texttt{/r0/sonar0} a \texttt{sonar7} & \texttt{sensor\_msgs/msg/Range}
  & Distancia media de cada sector.\\
\texttt{/r0/odom} & \texttt{nav\_msgs/msg/Odometry}
  & Posición y orientación para evaluación.\\
\texttt{/r0/collision} & \texttt{std\_msgs/msg/Bool}
  & Detección de colisión.\\
\bottomrule
\end{tabularx}
\end{center}

Los controladores producirán las salidas $V$ y $W$:
\centerEq{0\leq V\leq50\ \mathrm{km/h},
  \qquad -45^\circ\leq W\leq45^\circ.}
$W$ es el \textbf{mando de dirección en grados}, positivo hacia la izquierda.
El bloque de conversión proporcionado genera \texttt{linear.x} en
$\mathrm{m/s}$ y \texttt{angular.z} en $\mathrm{rad/s}$.
Se conservará este bloque en las dos partes de la práctica.

\newpage
\problem{4. Preparación y ejecución}

Se requiere MATLAB \textbf{R2025a o posterior}, con Simulink,
ROS Toolbox y Fuzzy Logic Toolbox. Para el controlador neuronal
se utilizarán las herramientas de aprendizaje estudiadas en la asignatura.

\romanitem{i}{Instalar el entorno siguiendo la guía del proyecto.
Desde la carpeta del repositorio, construir la imagen con
\texttt{make build} y lanzar la simulación con \texttt{make run}.
Abrir el
\href{http://localhost:6080/vnc.html?autoconnect=true&resize=scale}{simulador en el navegador}.}

\romanitem{ii}{Elegir un \texttt{ROS\_DOMAIN\_ID} y utilizar el mismo
valor en \texttt{docker/.env}, en \texttt{controller\_ros2.m} y en
\textit{ROS 2 / Configure ROS Network} del modelo Simulink.}

\romanitem{iii}{Trabajar en la carpeta \texttt{matlab/}.
El archivo \texttt{ackerman\_ROS2\_controller.slx} contiene la conexión
al robot, el controlador, la conversión de consignas y los bloques de
parada y registro. El controlador borroso deberá guardarse como
\texttt{fuzzy\_controller.fis}.}

\romanitem{iv}{Ejecutar \texttt{controller\_ros2} desde MATLAB.
Al terminar, \texttt{evaluate\_controller} presenta una figura con
el recorrido y los resultados.}

Para repetir una prueba desde la posición inicial, detener el simulador
con \textbf{Ctrl+C} y ejecutar de nuevo \texttt{make run}.
Desde otra terminal, \texttt{make attach} permite acceder al contenedor.
Si transcurre un segundo sin recibir órdenes, el coche del alumno se detiene.

\problem{5. Parte I. Controlador borroso Mamdani}

Se diseñará un controlador que utilice exclusivamente las medidas de los
sonares para regular la velocidad y el mando de dirección.
La posición y la orientación se reservan para la evaluación de la maniobra.

\romanitem{i}{Analizar la distribución de sensores y establecer qué
información permite reconocer rectas, curvas, proximidad al borde de la
carretera y vehículos que deben adelantarse.}

\romanitem{ii}{Diseñar un FIS de tipo Mamdani con ocho entradas, ordenadas
como \texttt{sonar\_0}, \ldots, \texttt{sonar\_7}, y dos salidas,
ordenadas como $V$ y $W$. El modelo mantiene las ocho entradas;
se justificará cuáles intervienen en las reglas de control.}

\romanitem{iii}{Definir y justificar los universos de discurso, las
funciones de pertenencia, la base de reglas, los operadores de inferencia
y el método de desborrosificación. Las distancias de entrada se
expresan en metros, con alcance máximo de $30\ \mathrm{m}$.
Antes del FIS, el modelo sustituye la ausencia de retorno por
$30\ \mathrm{m}$ y las lecturas inválidas por $0\ \mathrm{m}$.}

\romanitem{iv}{Integrar el controlador y ajustar su comportamiento mediante
pruebas sin tráfico y con tráfico. El diseño deberá permitir completar
el recorrido, adelantar y regresar al carril derecho.}

\romanitem{v}{Conservar los resultados de las pruebas y seleccionar una
configuración del controlador para generar los datos de entrenamiento
de la segunda parte.}

\newpage
\problem{6. Parte II. Controlador neuronal con LiDAR}

Se diseñará una red neuronal de regresión que reciba información del LiDAR
y genere las mismas consignas $V$ y $W$ que el controlador borroso.
Durante su ejecución, las entradas de la red procederán exclusivamente del
barrido LiDAR, o de un preprocesamiento de este.

\romanitem{i}{\textbf{Obtener los datos automáticamente.}
Ejecutar el controlador borroso y registrar, en cada muestra de control,
el barrido LiDAR y las consignas producidas por el borroso antes de la
conversión a \texttt{Twist}. Repetir recorridos que incluyan rectas,
curvas, aproximación al tráfico, adelantamiento y retorno al carril derecho.}

Para cada instante $k$, el conjunto de datos contendrá un par
\centerEq{\mathbf r_k=(r_{k,1},\ldots,r_{k,1081}),
  \qquad
  \mathbf y_k=
  \begin{bmatrix}V_k\\W_k\end{bmatrix}.}
Cada barrido se asociará con la consigna correspondiente a esa muestra,
conservando las marcas de tiempo y la identificación del recorrido.
Las paradas impuestas por fin de vuelta se tratarán separadamente del
comportamiento de conducción que se desea aprender.

\textbf{El archivo \texttt{data.mat} no contiene los barridos LiDAR.}
Este archivo registra la información de evaluación del recorrido.
Para entrenar la red será necesario añadir el registro sincronizado de
las medidas sensoriales y las consignas del controlador borroso.

\romanitem{ii}{\textbf{Preparar los datos.}
Definir un procedimiento reproducible para tratar medidas sin retorno
o inválidas, normalizar las entradas y preparar las salidas.
Puede utilizarse el barrido completo o una reducción de dimensión
justificada, obtenida directamente del LiDAR.
Los sonares, la posición y la orientación no se incorporarán como entradas
adicionales de la red.}

El preprocesamiento deberá proporcionar un vector de tamaño fijo y valores
finitos. El alcance máximo del sensor es $30\ \mathrm{m}$:
\texttt{+Inf} indica ausencia de retorno y \texttt{NaN} indica una
medida no válida. El mismo tratamiento se aplicará al entrenamiento
y a la ejecución en Simulink.

\romanitem{iii}{\textbf{Separar los datos.}
Crear conjuntos de entrenamiento, validación y prueba utilizando
recorridos distintos. Las muestras contiguas de un mismo recorrido
no se repartirán aleatoriamente entre entrenamiento y prueba.
Los parámetros de normalización que se estimen a partir de datos
se obtendrán únicamente del conjunto de entrenamiento.}

\romanitem{iv}{\textbf{Diseñar y entrenar la red.}
Seleccionar y justificar la arquitectura, las funciones de activación,
el procedimiento de entrenamiento y el criterio de selección del modelo.
Documentar los errores de predicción de $V$ y $W$ por separado y
comparar el comportamiento en entrenamiento, validación y prueba.}

\romanitem{v}{\textbf{Integrar y comprobar el control.}
Sustituir el bloque borroso por la red y su preprocesamiento.
Mantener el periodo $T_s=0{,}05\ \mathrm{s}$, la conversión de consignas,
los límites de salida y los bloques de parada y evaluación.
Comprobar la conducción en lazo cerrado y comparar sus resultados
con los del controlador borroso.}

\newpage
\problem{7. Evaluación y entrega}

Ambos controladores se evaluarán, como mínimo, en un escenario
\textbf{sin tráfico} y en otro \textbf{con tráfico}.
Las comparaciones se realizarán con la misma configuración del circuito,
del vehículo, de los sensores, del tráfico y de la posición inicial.
La configuración utilizada deberá quedar identificada en los resultados.

El modelo termina al completar una vuelta o detectar una colisión.
El estado de \texttt{/r0/collision} permanece activado hasta reiniciar
el simulador. El modelo proporcionado incluye además un tiempo máximo
de simulación de \textbf{$200\ \mathrm{s}$}.
Una ejecución que termine sin completar una vuelta se identificará
como \textbf{incompleta}.

Para cada ensayo se presentarán:
\begin{itemize}
  \setlength{\itemsep}{0pt}
  \item Vuelta completada o porcentaje recorrido y existencia de colisión.
  \item Tiempo de ejecución simulada.
  \item Tiempo y porcentaje de ocupación del carril contrario.
  \item Trayectoria y evolución de las consignas de velocidad y dirección.
\end{itemize}

El evaluador calcula la ocupación del carril contrario a partir de la
posición de referencia del vehículo en la odometría, situada en el
\textbf{eje trasero}. Las estadísticas de velocidad y dirección mostradas
en la figura corresponden a las \textbf{consignas del controlador}.
La detección de colisión se obtiene del topic correspondiente.

La memoria incluirá una comparación entre ambos controladores en los
dos escenarios. Se analizarán las maniobras de adelantamiento, los
fallos observados y la relación entre tiempo de vuelta y permanencia
en el carril contrario. En el controlador neuronal se distinguirá el
error de predicción sobre datos de prueba del comportamiento real
obtenido al conducir en lazo cerrado.

\begin{center}
\begin{tabularx}{0.96\textwidth}{@{}X r@{}}
\toprule
\textbf{Apartado} & \textbf{Peso}\\
\midrule
Control borroso: circuito sin tráfico & $25\,\%$\\
Control borroso: circuito con tráfico & $25\,\%$\\
Control neuronal: circuito sin tráfico & $15\,\%$\\
Control neuronal: circuito con tráfico & $15\,\%$\\
Memoria y presentación & $20\,\%$\\
\midrule
\textbf{Total} & \textbf{100\,\%}\\
\bottomrule
\end{tabularx}
\end{center}

Se podrá obtener hasta un \textbf{$10\,\%$ adicional por competición}.
Cada equipo seleccionará uno de sus controladores para ejecutarlo
en el escenario común establecido por el profesorado.
La clasificación tendrá en cuenta el tiempo de vuelta y la ocupación
del carril contrario, según los criterios anunciados antes de la prueba.
Se aportará una grabación continua del recorrido, sin cortes ni edición,
y la figura final de evaluación.

\romanitem{i}{\textbf{Código y modelos.}
Entregar el FIS, el controlador neuronal entrenado, los modelos Simulink
y los archivos necesarios para registrar datos, preprocesar, entrenar
y ejecutar ambos controladores.}

\romanitem{ii}{\textbf{Datos y reproducibilidad.}
Incluir los datos utilizados, o el procedimiento de acceso indicado
en el Aula Virtual, junto con la separación de recorridos, los parámetros
de entrenamiento y las configuraciones de prueba.}

\romanitem{iii}{\textbf{Memoria y presentación.}
Describir el diseño borroso, el conjunto de datos, el entrenamiento de la red
y la comparación experimental. Aportar las figuras de evaluación
y una demostración del funcionamiento de las dos propuestas.}

La fecha, el formato de entrega y las condiciones de la presentación
se publicarán en el Aula Virtual de la asignatura.

\end{document}
